% GENERATED by scripts/build.py convert ch07 -- do not hand-edit; edit
% chapters/md/07_conclusion.md and regenerate.
% md sha256 3db7ff6450af796cb344e0d9439bc00a183e010ae8462339e62f7405e7ec3d04
% canonical manuscript/v13_2026-10-17/07_conclusion.md sha256 3db7ff6450af796cb344e0d9439bc00a183e010ae8462339e62f7405e7ec3d04
% md identical to canonical: True
% filter sha256 d487b0e18e456f08097065d7ceaaf39fef0daee29958941b5539c501e0e4c24a

\section{Conclusion}\label{sec:7}\ThesisCheckNumber{section}{7}

This thesis examined the association between water and sanitation conditions and child anthropometric outcomes in four Sub-Saharan African countries: Ethiopia, Ghana, Kenya and Nigeria. Two distinctions ran throughout: household facilities versus surrounding community WASH, and a single survey-date exposure versus a preliminary, geospatially matched measure of early-life WASH coverage, built for an earlier survey round in three of the four countries, with Ethiopia retained only as a provisional sensitivity analysis (Section 1.3).

The baseline evidence shows no clear, robust association between improved household WASH and child outcomes within the same DHS cluster (Section 5.1): inconclusive, not evidence of no association. The community measure shows a large, significant association before covariate adjustment, attenuating to non-significance once socioeconomic status and region fixed effects are added (Section 5.2), consistent with, though not proof of, confounding by local development.

The historical extension's headline cell, the pooled established-date water--HAZ association (Section 5.6), is positive and significant; it carries the qualifications set out in full in Section 6 and cannot be read as evidence of an exposure-timing effect (Section 6.2). Neither analysis establishes a causal effect of WASH facilities on child outcomes, or a causal policy effect: the associations reported inform research and interpretation, not policy. The remaining gaps are detailed in Section 6 and Appendix C.
